% supplement_negative_results.tex
% Standalone paper: systematic closure of four minimal first-principles
% routes to dark energy in Einstein-Cartan-Holst gravity.
%
% Houston Golden, 2026
%
% Submitted to: [TBD — see publication-readiness checklist]
% arXiv: gr-qc / hep-th

\documentclass[11pt,letterpaper]{article}

\usepackage[utf8]{inputenc}
\usepackage[T1]{fontenc}
\usepackage{amsmath,amssymb,amsfonts}
\usepackage{mathtools}
\usepackage{physics}
\usepackage{geometry}
\geometry{margin=1in}
\usepackage{hyperref}
\usepackage{cleveref}
\usepackage{graphicx}
\usepackage{booktabs}
\usepackage{longtable}
\usepackage{enumitem}
\usepackage{slashed}

% Shortcuts
\newcommand{\MPl}{M_{\rm Pl}}
\newcommand{\Neff}{\Delta N_{\rm eff}}
\newcommand{\BI}{\gamma}               % Barbero-Immirzi
\newcommand{\Gsp}{G_{\rm SP}}
\newcommand{\Gcrit}{G_{\rm crit}}
\newcommand{\Geff}{G_{\rm eff}}
\newcommand{\Veff}{V_{\rm eff}}
\newcommand{\Gam}{\Gamma_{\rm eff}}
\newcommand{\Sfour}{S_{\rm 4f}}

\title{%
  Systematic Closure of Minimal First-Principles Routes\\
  to Dark Energy in Einstein-Cartan-Holst Gravity%
}
\author{Houston Golden\\
  Independent Researcher, Los Angeles, California, USA\\
  \texttt{houston@hubify.com}
}
\date{March 2026}

\begin{document}
\maketitle

% ------------------------------------------------------------------
\begin{abstract}
Einstein-Cartan-Holst (ECH) gravity---the Palatini formulation of general
relativity supplemented by the Holst term with Barbero-Immirzi parameter
$\BI$---is a natural candidate framework for geometric dark energy. The
minimal ECH$+$Dirac model generates a parity-odd vacuum operator through
spin-torsion coupling, and a phenomenological scaling argument relates
its inflationary dilution to the observed dark-energy density. However,
the late-time equation of state $w=-1$ must be assumed, not derived: it
requires showing that integrating out spin degrees of freedom generates
an IR-constant term in the effective action. We report four systematic
attempts to close this gap---two computational and two structural---together
with their negative outcomes:

\begin{enumerate}[label=(\roman*)]
  \item \textbf{Track~B (condensate route):} Fierz rearrangement of the
    torsion-induced four-fermion interaction shows that the
    scalar/pseudoscalar channel coupling
    $\Gsp\propto(\BI^2-1)/(\BI^2+1)$ is repulsive at $\BI<1$.
    Even when attractive ($\BI>1$), the coupling is ${\sim}175\times$ below
    the NJL critical coupling, and gravitational catalysis is exponentially
    suppressed. Gate~1 (condensate existence) fails; Track~B closes.
  \item \textbf{Branch~G v1 (strict one-loop effective action):} After exact
    torsion elimination, all finite-$\BI$ dependence resides in the
    four-fermion term $\Sfour$, which is quartic in fermion fields and does
    not contribute to the strict one-loop fermion determinant. The Dirac
    operator is the standard Levi-Civita operator with $\BI$-independent
    Seeley-DeWitt coefficients. Branch~G~v1 closes at Phase~1 (no novel
    Holst content at this order).
  \item \textbf{Route~T1 (dynamical Immirzi field):} Promoting the
    Barbero-Immirzi parameter $\BI\to\theta(x)$ to a dynamical pseudoscalar
    field. Torsion remains algebraic, but after torsion elimination the
    reduced action is identical to GR~$+$~standard axion-like particle (ALP).
    The four-fermion interaction is $\theta$-independent; no mechanism
    generates a potential $V(\theta)\sim\Lambda$; cosmological dynamics gives
    $w=+1$ (stiff matter), not $w=-1$. Route~T1 closes: the dynamical
    Immirzi field has no novel gravitational content beyond a standard ALP.
  \item \textbf{Route~S1 (parity-sensitive CMB phenomenology):} Assessment of
    whether the framework's parity-odd operator maps to distinctive
    birefringence observables. The minimal ECH$+$Dirac model has no
    photon-polarization coupling; any assumed coupling produces constant-$\beta$
    birefringence identical to generic ALP phenomenology with two free
    parameters for one observable. Route~S1 closes: no distinctive
    parity signal exists in the minimal model.
\end{enumerate}

These closures share a common structural origin: in the minimal ECH$+$Dirac
model, torsion is algebraic (non-propagating) and its elimination removes all
Holst-sector content from the reduced action. The Nieh-Yan topological
density carries no local dynamical information that survives this elimination.
The resulting physics is standard general relativity plus fermions plus a
constrained four-fermion contact interaction---insufficient for late-time
vacuum energy or distinctive observational signatures. The phenomenological
framework~\cite{Golden2026} survives as a motivated ansatz; progress toward
a first-principles derivation requires physics beyond the minimal model class.
All computations are symbolic or semi-analytic (Python/SymPy) and are fully
reproducible.
\end{abstract}

% ------------------------------------------------------------------
\section{Introduction}\label{sec:intro}

The cosmological constant problem---why $\rho_\Lambda \sim (2.3\;\text{meV})^4
\sim 10^{-122}\,\MPl^4$---motivates the search for geometric mechanisms that
generate or explain the observed vacuum energy from the structure of
spacetime itself. Einstein-Cartan-Holst (ECH) gravity is a natural candidate:
its Holst term introduces a parity-odd coupling controlled by the
Barbero-Immirzi parameter $\BI$, well-motivated by loop quantum
gravity~\cite{Freidel2005,Mercuri2009}. The spin-torsion interaction
generates four-fermion contact terms and a parity-violating vacuum
operator, suggesting a possible geometric origin for dark energy.

A phenomenological framework based on this idea~\cite{Golden2026} constructs
a dark-energy density $\rho_\Lambda = \Xi\,\MPl^4$ through inflationary
dilution of the parity-odd operator, reducing fine-tuning from $10^{120}$ to
${\sim}10^5$. However, the framework contains a logical gap: the operator
$\langle\varepsilon^{abcd}K_{ab}R_{cd}\rangle$ is sourced by the fermion
spin density, which vanishes at late times. For the residual to persist as a
genuine vacuum energy with $w=-1$, one must show that integrating out the spin
degrees of freedom generates an IR-constant term in the effective action.

This paper reports a structured program to close that gap. Four routes were
tested---two computational, one structural extension, and one
phenomenological assessment. All four fail. These are clean negative results:
they narrow the landscape of viable derivation strategies, close the most
natural minimal-model candidates, and establish concrete structural
requirements for any next-generation geometric dark-energy mechanism.

\subsection*{Scope and limitations}
This paper reports results for the \emph{minimal} ECH$+$Dirac model
(Einstein-Cartan gravity with Holst term, minimally coupled to $N_f$ Dirac
fermions, at the LQG Barbero-Immirzi value $\BI=0.274$), together with the
simplest structural extension (promoting $\BI$ to a dynamical field). Routes
involving non-minimal couplings, higher-loop effects, non-perturbative
dynamics, graviton-loop contributions, genuinely new symmetry structures,
or different values of~$\BI$ are not tested and remain open. The negative
results reported here apply only at the stated approximation orders and
model classes.


% ------------------------------------------------------------------
\section{Context and Motivation}\label{sec:context}

The phenomenological framework of Ref.~\cite{Golden2026} treats $w=-1$ as an
explicit assumption, stating: ``We therefore treat $w=-1$ as a
phenomenological assumption [\ldots] and flag the derivation of the late-time
IR effective action as an important open problem.'' The present paper
constitutes a systematic test of whether the minimal microscopic theory
delivers such a term.

The question is not merely whether a specific calculation succeeds, but
whether the \emph{model class}---minimal ECH gravity with algebraic torsion
and standard Dirac coupling---has the structural capacity to generate
late-time vacuum physics. The four routes tested here cover the most natural
candidates within this model class:
\begin{itemize}
  \item condensation of the torsion-induced four-fermion interaction,
  \item the one-loop fermion effective action,
  \item promotion of $\BI$ to a dynamical field, and
  \item mapping the parity-odd operator to distinctive observables.
\end{itemize}
The simultaneous failure of all four points to a structural limitation of the
model class itself, not merely to the failure of individual calculations.

% ------------------------------------------------------------------
\section{Minimal Model and Tested Question}\label{sec:model}

\subsection{Microscopic action}

The starting point is the Einstein-Cartan-Holst action minimally coupled to
$N_f$ Dirac fermions:
\begin{align}
  S &= S_{\rm grav}[e,\omega] + S_{\rm ferm}[e,\omega,\psi], \label{eq:action}\\
  S_{\rm grav} &= \frac{\MPl^2}{2}\int
    \bigl[\varepsilon_{IJKL} + \tfrac{2}{\BI}\,\eta_{I[K}\eta_{L]J}\bigr]\,
    e^I\wedge e^J\wedge F^{KL}[\omega], \label{eq:Sgrav}\\
  S_{\rm ferm} &= \int d^4x\,|e|\,
    \bigl[\bar\psi_i\,i\gamma^\mu D_\mu\psi_i - m_i\bar\psi_i\psi_i\bigr],
    \label{eq:Sferm}
\end{align}
with $D_\mu = \partial_\mu + \tfrac{1}{4}\omega_\mu^{IJ}\gamma_I\gamma_J$,
metric signature $(+,-,-,-)$, $\gamma^5=i\gamma^0\gamma^1\gamma^2\gamma^3$
with $(\gamma^5)^2=1$, and $\kappa^2 = 8\pi G = \MPl^{-2}$. The
Barbero-Immirzi parameter~$\BI$ is a fixed constant. No bare cosmological
constant is included in the starting action.

\subsection{The central question}

Does integrating out fermion fields in the action~\eqref{eq:action} generate a
late-time vacuum term with $w\approx -1$ that is not merely a restatement of
bare-$\Lambda$ renormalization and that depends nontrivially on finite~$\BI$?

Four routes were tested, corresponding to different approximation strategies
and extensions:
\begin{itemize}
  \item \textbf{Track~B:} Test whether a pseudoscalar fermion condensate
    $\Phi = \langle\bar\psi i\gamma^5\psi\rangle$ forms via the NJL mechanism
    from the torsion-induced four-fermion interaction and persists as
    late-time vacuum energy.
  \item \textbf{Branch~G v1:} Compute the strict one-loop gravitational
    effective action $\Gam[e]$ after torsion elimination and fermion
    integration, and test whether it contains a $\BI$-dependent vacuum-like
    term beyond standard cosmological-constant renormalization.
  \item \textbf{Route~T1:} Promote the Barbero-Immirzi parameter to a
    dynamical pseudoscalar field $\BI\to\theta(x)$ and assess whether the
    resulting theory contains novel gravitational content (non-ALP terms,
    $\theta$-dependent four-fermion structure, or a dynamically generated
    vacuum potential).
  \item \textbf{Route~S1:} Assess whether the framework's parity-odd operator
    $\varepsilon^{abcd}K_{ab}R_{cd}$ maps to distinctive, testable
    birefringence observables that are not reproducible by generic
    axion-like-particle models.
\end{itemize}


% ------------------------------------------------------------------
\section{Track~B: Condensate Route}\label{sec:trackB}

\subsection{Setup}\label{sec:trackB_setup}

The spin connection $\omega^{IJ}_\mu$ is non-propagating in EC gravity; its
equation of motion is algebraic. Solving it exactly yields the reduced
action~\cite{Freidel2005,ShapiroTeixeira2014}
\begin{equation}\label{eq:Sreduced}
  S_{\rm reduced}[e,\psi] = S_{\rm EH}[e]
  + \int d^4x\sqrt{-g}\,\bigl[\bar\psi\,i\gamma^\mu\mathring\nabla_\mu\psi
    - m\bar\psi\psi\bigr]
  + \Sfour[e,\psi;\BI],
\end{equation}
where $\mathring\nabla$ uses the torsion-free Levi-Civita connection and
$\Sfour$ is the torsion-induced four-fermion interaction.

\paragraph{Perfect-square structure.}
The four-fermion term has a constrained form discovered during
Computation~1~\cite{comp1}:
\begin{equation}\label{eq:Sfour}
  \mathcal{L}_{\rm 4f}
  = -\Geff\bigl(\bar\psi\gamma^\mu\gamma^5\psi
    + \tfrac{1}{\BI}\,\bar\psi\gamma^\mu\psi\bigr)^2,
  \qquad
  \Geff = \frac{3\kappa^2}{16}\,\frac{\BI^2}{\BI^2+1}.
\end{equation}
This is a perfect square of a single combined current
$J^\mu = J^\mu_{(A)} + (1/\BI)\,J^\mu_{(V)}$. The individual couplings
$G_V$, $G_A$, $G_{VA}$ are \emph{not} independent parameters; they are
constrained by this squared-current structure.

At $\BI=0.274$, the vector admixture $1/\BI\approx 3.65$ is
$\mathcal{O}(1)$---the vector sector dominates over the axial sector.

\subsection{Gate structure}\label{sec:trackB_gates}

Track~B was organized around three quantitative gates, defined and frozen
before computation began~\cite{canonical_v3}:
\begin{description}[style=nextline]
  \item[Gate~1 (Condensate existence)]
    The renormalized effective potential $\Veff(\sigma,\pi)$ admits a local
    minimum at $(\sigma^*,\pi^*)$ with $\pi^*\neq 0$ and positive-definite
    Hessian.
  \item[Gate~2 (Vacuum persistence)]
    The minimum persists when curvature $R\to 0$ and spin source $s\to 0$.
  \item[Gate~3 (Vacuum-like stress-energy)]
    The condensate contribution yields $\rho_\Phi > 0$ and $|1+w|<10^{-2}$,
    with scheme-independent existence/sign.
\end{description}
Gate~1 failed; Gates~2 and~3 were not reached.

\subsection{Algebraic and numerical results}\label{sec:trackB_results}

\paragraph{Result~1: The scalar/pseudoscalar channel is repulsive at
$\BI=0.274$.}
Fierz rearrangement of the four-fermion interaction~\eqref{eq:Sfour}
decomposes it into standard NJL channels~\cite{comp2,comp2b}. The effective
coupling in the scalar/pseudoscalar sector is
\begin{equation}\label{eq:GSP}
  \Gsp = \frac{3\kappa^2}{16}\,\frac{\BI^2-1}{\BI^2+1}.
\end{equation}
This is an exact algebraic result. The sign changes at $\BI=1$:
\begin{itemize}
  \item $\BI>1$: attractive (condensation possible in principle),
  \item $\BI=1$: decoupled,
  \item $\BI<1$: \textbf{repulsive} (condensation impossible via NJL).
\end{itemize}
At $\BI=0.274<1$, the scalar/pseudoscalar channel is repulsive. The NJL gap
equation has no nontrivial solution for any repulsive coupling, regardless of
curvature or fermion content.

\paragraph{Result~2: Even with the correct sign, the coupling is
${\sim}175\times$ below NJL criticality.}
The NJL critical coupling for condensation in proper-time regularization is
$\Gcrit = 2\pi^2/(N_c\Lambda^2)$~\cite{comp3b}. With gravitational coupling
$G\sim\kappa^2\sim 1/\MPl^2$ and cutoff $\Lambda\sim\MPl$:
\begin{equation}\label{eq:ratio}
  \frac{\Gsp}{\Gcrit}\bigg|_{\BI=2.0}\approx 0.006.
\end{equation}
The gravitational four-fermion coupling is subcritical by two orders of
magnitude, even when the sign is attractive.

\paragraph{Result~3: Gravitational catalysis is exponentially suppressed.}
Gravitational catalysis~\cite{GorbarGusynin2008} lowers the threshold for
\emph{already-attractive} channels but cannot flip the sign of a repulsive
channel. For the attractive case ($\BI>1$) at Planck curvature
$R\sim\MPl^2$~\cite{comp3b}:
\begin{equation}\label{eq:catalysis}
  M_*\sim\Lambda\exp\!\bigl(-\text{const}/(G_{\rm SP}\times R)\bigr)
  \sim\exp(-2100)\approx 10^{-914}\,\MPl.
\end{equation}
The corresponding vacuum energy $V_{\rm vac}\sim M_*^4\sim 10^{-3658}\,\MPl^4$
is smaller than the observed $\rho_\Lambda\sim 10^{-122}\,\MPl^4$ by
${\sim}3500$ orders of magnitude.

\paragraph{Result~4: Multi-flavor and alternative channels do not rescue the
mechanism.}
\begin{itemize}
  \item \emph{Multi-flavor:} Inter-flavor contributions to the S/P channel
    vanish because $\tr(\gamma^\mu)=0$ and $\tr(\gamma^\mu\gamma^5)=0$; the
    Fierz traces factorize for different species~\cite{comp3}.
  \item \emph{Vector/axial condensation:} The VV and AA channels \emph{are}
    attractive, but a vector condensate breaks Lorentz invariance and does not
    produce isotropic vacuum energy compatible with dark energy.
  \item \emph{VA cross-term:} The VA Fierz coefficient vanishes identically
    for identical single-flavor fermions (Fierz exchange symmetry). The
    effective potential has $O(2)$ symmetry in $(\sigma,\pi)$; Possibility~B
    from the canonical problem statement is realized, not
    Possibility~A~\cite{comp2b}.
\end{itemize}

\subsection{Closure of Track~B}\label{sec:trackB_closure}

Gate~1 fails on four independent lines of evidence: wrong sign, subcritical
coupling, exponentially suppressed catalysis, and no multi-flavor rescue.
Track~B closes for the minimal single-flavor NJL pseudoscalar condensate at
$\BI=0.274$.

\Cref{fig:GSP} shows $\Gsp$ as a function of~$\BI$, with the sign flip at
$\BI=1$ and the LQG value $\BI=0.274$ marked. \Cref{fig:veff} compares
$\Veff$ in the repulsive ($\BI=0.274$) and attractive ($\BI=2.0$) regimes.


\begin{figure}[t]
  \centering
  \includegraphics[width=0.75\textwidth]{figures/G_SP_vs_gamma.png}
  \caption{Effective scalar/pseudoscalar coupling $\Gsp$ (Eq.~\ref{eq:GSP})
    as a function of the Barbero-Immirzi parameter~$\BI$. The coupling
    changes sign at $\BI=1$ (vertical dashed line). At the LQG value
    $\BI=0.274$ (red marker), $\Gsp<0$: the channel is repulsive and NJL
    condensation is impossible.}
  \label{fig:GSP}
\end{figure}

\begin{figure}[t]
  \centering
  \includegraphics[width=0.75\textwidth]{figures/veff_repulsive_vs_attractive.png}
  \caption{One-loop effective potential $\Veff(\pi)$ at fixed $\sigma=0$,
    comparing the repulsive regime ($\BI=0.274$, no nontrivial minimum) with
    the attractive regime ($\BI=2.0$). Even in the attractive case, the
    coupling is ${\sim}175\times$ below the NJL critical coupling, so no
    condensation occurs.}
  \label{fig:veff}
\end{figure}


% ------------------------------------------------------------------
\section{Branch~G v1: Strict One-Loop Effective-Action Route}
\label{sec:branchG}

\subsection{Setup}\label{sec:branchG_setup}

Branch~G v1 asks a different question from Track~B: does the one-loop
gravitational effective action $\Gam[e]$---obtained by integrating out all
fermions on a curved background---contain a $\BI$-dependent vacuum-like term
that goes beyond standard cosmological-constant renormalization?

The target object is
\begin{equation}\label{eq:Gam}
  \Gam[e] = S_{\rm grav}[e,\mathring\omega[e]]
  + \Gamma^{(1\text{-loop})}_{\rm ferm}[e],
\end{equation}
where $\Gamma^{(1\text{-loop})}_{\rm ferm}
= -\tfrac{1}{2}\,\mathrm{Tr}\ln(-\slashed{D}^2 + \ldots)$ is the one-loop
fermion determinant.

\paragraph{Scope restriction.}
Branch~G v1 is restricted to \emph{Approach~B} (strict one-loop): the
four-fermion term $\Sfour$ is treated as a vertex in the loop expansion and
does not contribute at one loop. The Hubbard-Stratonovich route
(Approach~A), which would reintroduce condensate-like auxiliary-field
structure, is out of scope and would require a separate branch
definition~\cite{branchG_canonical}.

\subsection{Gate structure}\label{sec:branchG_gates}

Three gates were defined and frozen before computation~\cite{branchG_canonical}:
\begin{description}[style=nextline]
  \item[Gate~G1 (Nontrivial finite-$\BI$ vacuum structure)]
    The renormalized $\Gam[e]$ contains a cosmological-constant-like term
    whose existence depends nontrivially on finite~$\BI$ and is not merely
    standard bare-$\Lambda$ renormalization.
  \item[Gate~G2 (IR physicality)]
    The term survives as $R\to 0$ and is scheme-robust in existence/sign.
  \item[Gate~G3 (Dark energy viability)]
    Positive vacuum energy, $|1+w|<10^{-2}$, mappable to the main paper's
    dark-energy sector.
\end{description}
An additional early-kill condition (FM-G6) was predeclared: if the strict
one-loop object has no novel finite-$\BI$ content whatsoever, Branch~G v1
closes at Phase~1 without reaching Gates~G1--G3.

\subsection{Structural argument}\label{sec:branchG_argument}

After exact torsion elimination, the reduced action is
Eq.~\eqref{eq:Sreduced}. The free Dirac term uses the Levi-Civita connection
$\mathring\nabla$. All $\BI$-dependence resides in $\Sfour$, which is quartic
in~$\psi$.

\paragraph{The one-loop operator is $\BI$-independent.}
At strict one-loop (Approach~B), the Dirac operator whose determinant is
computed is
\begin{equation}\label{eq:Dirac}
  \slashed{D} = i\gamma^\mu\mathring\nabla_\mu,
\end{equation}
the standard curved-space Dirac operator with Levi-Civita connection. This
operator does not depend on~$\BI$. The four-fermion term $\Sfour$ is quartic
in~$\psi$ and contributes only at two loops or beyond.

Therefore, $\Gamma^{(1\text{-loop})}_{\rm ferm}[e]$ is $\BI$-independent as a
functional. Its Seeley-DeWitt coefficients $a_0$, $a_1$, $a_2,\ldots$ are the
standard textbook coefficients for a spin-$\tfrac{1}{2}$ field on a Riemannian
background~\cite{Vassilevich2003}, with no $\BI$-dependence.

\paragraph{Three candidate channels for $\BI$ entry.}
Before declaring closure, three channels through which $\BI$ might enter at
one loop were checked~\cite{branchG_analysis}:

\begin{enumerate}[label=(\alph*)]
  \item \textbf{Modified background equations:} The classical field equations
    include $\BI$-dependent terms, so the background geometry (e.g., the
    bounce solution) depends on~$\BI$. However, this is $\BI$-dependence of
    the \emph{field configuration} (a point in field space), not of the
    effective action \emph{as a functional}. This does not constitute novel
    Holst content.
  \item \textbf{Parity-odd sector / $\eta$-invariant:} The fermion determinant
    phase (Atiyah-Patodi-Singer $\eta$-invariant) is computed from the
    spectrum of~$\slashed{D}=i\gamma^\mu\mathring\nabla_\mu$, which is
    $\BI$-independent. Furthermore, the Pontryagin density $P_4=0$ on
    homogeneous FRW backgrounds (conformally flat), and even on backgrounds
    where $P_4\neq 0$, the contribution is curvature-local and topological,
    not vacuum-like.
  \item \textbf{Functional measure corrections:} Torsion elimination is a
    classical algebraic substitution, not a change of fermionic integration
    variable. The standard treatment~\cite{Freidel2005,Shapiro2002} performs
    this without measure corrections. The path integral over~$\omega$ is
    Gaussian in~$\omega$; the resulting determinant factor depends on the
    tetrad~$e$ but not on~$\psi$, modifying the gravitational measure rather
    than the fermion determinant.
\end{enumerate}

\paragraph{Literature confirmation.}
Two existing one-loop computations were examined:
\begin{itemize}
  \item Chattopadhyay~(2023)~\cite{Chattopadhyay2023} computes the
    \emph{graviton} (not fermion) one-loop effective action in chiral EC
    gravity. No fermions are included; $\BI$ does not appear; the computation
    is for a different object.
  \item Shapiro \& Teixeira~(2014)~\cite{ShapiroTeixeira2014} find
    $\BI$-dependent one-loop divergences, but these arise from the
    \emph{graviton} loop with \emph{external} fermion currents treated as
    classical backgrounds. Fermions are not quantum fields in their
    computation. This confirms that $\BI$ enters through graviton vertices,
    not the fermion operator.
\end{itemize}

\subsection{Closure of Branch~G v1}\label{sec:branchG_closure}

\paragraph{Novel Holst content criterion (from the frozen canonical
statement):}
\begin{quote}
``Branch~G counts as nontrivially alive only if, at the stated approximation
order, the renormalized effective action $\Gam[e]$ contains at least one
finite or regulator-robust term whose existence, sign, or classification
depends nontrivially on finite~$\BI$ and either disappears or changes
qualitatively in the $\BI\to\infty$ Einstein-Cartan limit.''
\end{quote}

No such term exists at strict one-loop. FM-G6 fires. Branch~G~v1 closes at
Phase~1 without proceeding to Gates~G1--G3.


% ------------------------------------------------------------------
\section{Route~T1: Dynamical Immirzi Field}\label{sec:T1}

\subsection{Motivation and setup}\label{sec:T1_setup}

The Barbero-Immirzi parameter $\BI$ appears in the Holst
action~\eqref{eq:Sgrav} as a dimensionless coupling. Promoting it to a
dynamical pseudoscalar field $\BI\to\theta(x)$ is a natural structural
extension: $\theta$ then plays the role of an axion-like field whose
``decay constant'' is tied to the gravitational sector through the Nieh-Yan
topological density~\cite{Mercuri2009}.

The question tested in Route~T1 is whether this promotion generates
genuinely new gravitational content---specifically, whether the reduced
action after torsion elimination contains terms that are \emph{not}
reproducible by simply adding a standard ALP to general relativity.

\subsection{Literature review results}\label{sec:T1_literature}

Six key references were examined~\cite{Mercuri2009,CalcagniMercuri2009,%
TaverasYunes2008,TorresGomez2009,LattanziMercuri2010,Magueijo2013}.
The assessment addressed six questions (Q1--Q6), each with predefined
kill criteria.

\paragraph{Q1: Torsion remains algebraic.}
Even with $\theta(x)$ dynamical, the torsion equation of motion remains
algebraic in the contorsion tensor $K^I{}_{JK}$. No derivatives of~$K$
appear. Torsion elimination proceeds exactly as in the constant-$\BI$
case~\cite{CalcagniMercuri2009}. This is a necessary (but not sufficient)
condition for tractability.

\paragraph{Q2: The reduced action is a standard ALP.}
After torsion elimination with dynamical $\theta(x)$, the reduced action
takes the form~\cite{CalcagniMercuri2009,Mercuri2009}
\begin{equation}\label{eq:T1reduced}
  S_{\rm reduced}[g,\psi,\theta] = S_{\rm EH}[g]
  + \int d^4x\sqrt{-g}\left[
    -\tfrac{1}{2}f(\theta)(\partial\theta)^2
    + \bar\psi\,i\gamma^\mu\mathring\nabla_\mu\psi
    - m\bar\psi\psi
    + \mathcal{L}_{\rm 4f}
  \right],
\end{equation}
where $f(\theta)$ is a kinetic function and $\mathcal{L}_{\rm 4f}$ is the
four-fermion interaction. The critical structural result:

\emph{Every term in $S_{\rm reduced}$ is of a standard form.}
The kinetic term $f(\theta)(\partial\theta)^2$ is standard ALP kinetic
energy. The Nieh-Yan coupling $\theta\,dT^I\wedge e_I$ reduces, after
torsion elimination, to a boundary term plus standard axion-type
couplings. No term arises that is not reproducible by adding a
conventional pseudoscalar field to GR.

\paragraph{Q2 (continued): Four-fermion interaction is $\theta$-independent.}
The torsion-induced four-fermion interaction after elimination is
\begin{equation}\label{eq:T1fourfermi}
  \mathcal{L}_{\rm 4f} = -\frac{3}{2}\pi G\,
  (\bar\psi\gamma^\mu\gamma^5\psi)(\bar\psi\gamma_\mu\gamma^5\psi).
\end{equation}
This is the standard axial-axial contact interaction and is
\emph{independent of $\theta$}. The Fierz channel structure is unchanged
from the constant-$\BI$ case analyzed in Track~B. In particular, promoting
$\BI\to\theta(x)$ does not modify the scalar/pseudoscalar channel coupling
or introduce new condensation channels. The Track~B negative results carry
over unchanged.

\paragraph{Q3: Literature is complete.}
The computation of the reduced action with dynamical Immirzi field has
been performed by Calcagni \& Mercuri~\cite{CalcagniMercuri2009} and
independently by Taveras \& Yunes~\cite{TaverasYunes2008}. The
cosmological dynamics have been analyzed by
Lattanzi \& Mercuri~\cite{LattanziMercuri2010}. No gap in the literature
was identified.

\paragraph{Q4: No mechanism generates $V(\theta)\sim\Lambda$.}
The Mercuri~\cite{Mercuri2009} analogy posits that $\theta$ plays the role
of the QCD axion, with the Nieh-Yan density replacing $F\wedge\tilde{F}$.
However:
\begin{itemize}
  \item In QCD, the axion potential is generated by instantons with a
    definite non-perturbative scale $\Lambda_{\rm QCD}$. The resulting
    axion mass is $m_a\sim\Lambda_{\rm QCD}^2/f_a\sim 10^{-12}$~eV.
  \item In gravity, the analogous scale would be set by gravitational
    instantons, which are Planck-suppressed: $V(\theta)\sim\MPl^4\,
    e^{-S_{\rm inst}}$ with $S_{\rm inst}\sim\MPl^2/H^2\gg 1$ at
    late times.
  \item The resulting vacuum energy is ${\sim}10^{-36}$ orders of
    magnitude below $\rho_\Lambda\sim 10^{-122}\,\MPl^4$.
  \item Any $V(\theta)$ yielding the observed $\Lambda$ must be put in
    by hand. This displaces the fine-tuning problem from ``why $w=-1$?''
    to ``why $V(\theta_0)=\Lambda_{\rm obs}$?''
\end{itemize}

\paragraph{Q5: Cosmological dynamics give $w=+1$, not $w=-1$.}
Taveras \& Yunes~\cite{TaverasYunes2008} show that the dynamical Immirzi
field on cosmological backgrounds behaves as stiff matter with equation of
state $w=+1$. This is the \emph{opposite} of dark energy. The field
redshifts as $\rho_\theta\propto a^{-6}$, becoming negligible at late
times. There is no mechanism for the field to drive late-time acceleration.

\paragraph{Q6: The Nieh-Yan topological nature is fatal.}
The Nieh-Yan 4-form $N_4 = d(e^I\wedge T_I)$ is exact (a total derivative)
in four dimensions. Its topological nature ensures that any coupling of
$\theta$ to $N_4$ reduces to boundary terms plus standard kinetic and
coupling structures after torsion elimination. This is the structural
reason why the gravitational origin washes out: the Nieh-Yan density
carries no local dynamical information that survives torsion
elimination~\cite{Mercuri2009,CalcagniMercuri2009}.

\subsection{Closure of Route~T1}\label{sec:T1_closure}

Route~T1 closes on three independent grounds:
\begin{enumerate}
  \item \textbf{No novel content:} The reduced action after torsion
    elimination is identical to GR~$+$~standard ALP. Every term is
    reproducible without the spin-torsion origin.
  \item \textbf{No vacuum mechanism:} No dynamical mechanism generates
    $V(\theta)\sim\Lambda$. Gravitational instantons are
    Planck-suppressed; any viable potential must be assumed, not derived.
  \item \textbf{Wrong cosmological dynamics:} The field behaves as stiff
    matter ($w=+1$), the opposite of dark energy.
\end{enumerate}

Route~T1 was assessed through literature review only; no original symbolic
computation was performed, as the existing literature fully covers the
relevant calculations.


% ------------------------------------------------------------------
\section{Route~S1: Parity-Sensitive CMB Phenomenology}\label{sec:S1}

\subsection{Motivation and setup}\label{sec:S1_setup}

The ECH action contains a parity-odd operator through the Holst term
$\varepsilon^{abcd}K_{ab}R_{cd}$, which sources the framework's
parity-violating vacuum structure. Cosmic birefringence---rotation of CMB
polarization by an angle $\beta$---has been detected at combined
significance $\beta = 0.342^\circ \pm 0.094^\circ$ (isotropic,
$3.6\sigma$)~\cite{EskiltKomatsu2022} with recent ACT DR6 results
strengthening the evidence~\cite{DiegoPalazuelos2025}. Route~S1 asks
whether the framework predicts a distinctive birefringence signal.

\subsection{Assessment results}\label{sec:S1_results}

Four questions (Q1--Q4) were assessed, each with predefined kill criteria.

\paragraph{Q1: The minimal framework has no photon coupling.}
The parity-odd operator $\varepsilon^{abcd}K_{ab}R_{cd}$ is purely
gravitational. It couples to the fermion spin density through torsion but
has \emph{no coupling to the photon sector} in the minimal ECH$+$Dirac
model. The electromagnetic field strength $F_{\mu\nu}$ does not appear in
the Holst term, the torsion equation of motion, or the reduced action
after torsion elimination.

For the framework to produce birefringence, one must \emph{assume} a
coupling of the form $\theta F\wedge\tilde{F}$ or equivalently
$g_{\gamma}\,\theta\,\varepsilon^{\mu\nu\rho\sigma}F_{\mu\nu}F_{\rho\sigma}$,
with $g_\gamma$ as a free parameter. This coupling is not derived from the
microscopic ECH action; it must be added by hand.

The main paper itself acknowledges this gap (Sec.~III.B.2 of
Ref.~\cite{Golden2026}): the birefringence discussion is framed as a
consistency check, not a prediction.

\paragraph{Q2: Parameterization is unfalsifiable.}
Even with an assumed photon coupling, the predicted birefringence angle is
\begin{equation}\label{eq:beta}
  \beta = f_{\rm photon}\,C_0\,\frac{\alpha}{M}\,\MPl,
\end{equation}
where $f_{\rm photon}$ is the assumed coupling strength (free),
$C_0$ encodes the dilution history (free within the framework's scaling
ansatz), $\alpha/M$ is the framework's dark-energy scale parameter. This
formula has two free parameters ($f_{\rm photon}$, $C_0$) for one
observable ($\beta$). Any measured $\beta$ can be accommodated by tuning
these parameters. The mapping is unfalsifiable.

\paragraph{Q3: No distinguishing power from generic ALP birefringence.}
The predicted signal is constant-$\beta$ (uniform rotation independent of
multipole $\ell$). This is identical to the prediction of a generic
axion-like particle with constant coupling. The spin-torsion origin
provides:
\begin{itemize}
  \item no distinctive $\beta(\ell)$ shape (scale-dependent rotation),
  \item no cross-correlation with other observables ($\Neff$, galaxy
    spins) that would distinguish it from a generic ALP,
  \item no predicted sign or magnitude constraint on $\beta$ from the
    framework's other fitted parameters.
\end{itemize}
All three survival criteria from the Phase~1 checklist fail.

\paragraph{Q4: Data exists but cannot be used distinctively.}
Current data (Planck NPIPE, ACT DR6) provides excellent sensitivity to
isotropic birefringence~\cite{DiegoPalazuelos2025,EskiltKomatsu2022}.
However, since the framework's prediction is indistinguishable from
generic constant-$\beta$ ALP birefringence, fitting the data would
produce a consistency check (``the framework is compatible with observed
$\beta$ for $f_{\rm photon}\sim\mathcal{O}(1)$''), not a test of the
spin-torsion mechanism.

\subsection{Closure of Route~S1}\label{sec:S1_closure}

Route~S1 closes on three independent grounds:
\begin{enumerate}
  \item \textbf{No photon coupling:} The minimal framework has no
    mechanism to rotate photon polarization. Any coupling must be
    assumed, not derived.
  \item \textbf{Unfalsifiable parameterization:} Two free parameters for
    one observable; any $\beta$ can be accommodated.
  \item \textbf{No distinguishing power:} The prediction is identical to
    generic ALP birefringence; no distinctive spin-torsion signature
    exists.
\end{enumerate}

Route~S1 was assessed through literature review and structural analysis;
no pipeline computation was performed, as the assessment determined that
no distinctive prediction exists to test.


% ------------------------------------------------------------------
\section{Combined Interpretation}\label{sec:combined}

\Cref{tab:summary} summarizes the four tested routes and their outcomes.

\begin{table}[t]
\centering
\caption{Summary of four tested routes to first-principles dark energy or
distinctive observables in the ECH framework, and their outcomes.}
\label{tab:summary}
\small
\begin{tabular}{@{}lllll@{}}
\toprule
Route & Target & Method & Outcome & Failure point \\
\midrule
Track~B & Condensate
  & Fierz + one-loop $\Veff$
  & Gate~1 failed
  & S/P channel repulsive \\
  & $\langle\bar\psi i\gamma^5\psi\rangle$
  & & & at $\BI=0.274$ \\[3pt]
Branch~G v1 & Effective action $\Gam[e]$
  & Strict one-loop
  & FM-G6 fired
  & Fermion determinant \\
  & & after torsion elim.
  &
  & is $\BI$-independent \\[3pt]
Route~T1 & Dynamical Immirzi
  & Literature review
  & All gates fail
  & Reduced action $=$  \\
  & field $\theta(x)$
  & + structural analysis
  &
  & standard ALP \\[3pt]
Route~S1 & Parity/birefringence
  & Structural assessment
  & Gate S1-1 failed
  & No photon coupling; \\
  & observables
  &
  &
  & $\beta=\,$generic ALP \\
\bottomrule
\end{tabular}
\end{table}

The minimal ECH$+$Dirac model, at the tested approximation orders and
extensions, does not generate the late-time vacuum term from first principles
via any of the four routes tested:
\begin{enumerate}[label=(\roman*)]
  \item the NJL condensate route (Track~B),
  \item the strict one-loop reduced effective-action route (Branch~G v1),
  \item the dynamical Immirzi field extension (Route~T1), or
  \item the parity-sensitive CMB phenomenology channel (Route~S1).
\end{enumerate}

\paragraph{Structural lesson.}
The four closures converge on a common structural explanation. In the
minimal ECH$+$Dirac model, torsion is algebraic (non-propagating). Solving
the torsion equation of motion exactly eliminates all Holst-sector
structure, leaving only:
\begin{itemize}
  \item the standard Einstein-Hilbert gravitational sector,
  \item standard minimally-coupled fermion kinetic terms, and
  \item a four-fermion contact interaction with constrained channel
    structure.
\end{itemize}
The Barbero-Immirzi parameter enters only the four-fermion coupling
constants, which are subcritical for condensation (Track~B) and invisible
at strict one-loop (Branch~G v1). Promoting $\BI\to\theta(x)$ does not
escape this structural trap: the Nieh-Yan density is topological (exact
in 4D), so the dynamical field reduces to a standard ALP after torsion
elimination (Route~T1). And the parity-odd operator, being purely
gravitational, has no photon coupling to produce distinctive birefringence
(Route~S1).

\paragraph{What this does not mean.}
These closures do \emph{not} invalidate the main phenomenological framework.
The main paper already treats $w=-1$ as an assumption, and these results
confirm that this treatment is appropriate. The negative results establish
the boundary of what the minimal model achieves at the lowest nontrivial
approximation orders.

\paragraph{What remains open.}
The following routes were not tested and remain logically open, each
requiring its own canonical problem statement:
\begin{itemize}
  \item \emph{Graviton one-loop with fermion currents}
    (Shapiro-Teixeira approach): $\BI$ enters through graviton vertices.
    This is a graviton-loop computation, structurally different from the
    fermion-loop computations tested here.
  \item \emph{Beyond-one-loop / resummation:} The four-fermion interaction
    $\Sfour$ contributes at two loops. Hubbard-Stratonovich resummation
    (Approach~A) may introduce $\BI$-dependent condensate channels, but
    Track~B showed the S/P channel fails even at this level.
  \item \emph{Non-perturbative effects:} Instantons, tunneling, topological
    contributions.
  \item \emph{Non-minimal couplings:} Additional torsion-matter interactions
    beyond the minimal EC$+$Holst$+$Dirac sector.
  \item \emph{New symmetry structures:} Genuinely new light degrees of
    freedom with symmetry-protected masses, beyond the minimal model
    class tested here.
\end{itemize}


% ------------------------------------------------------------------
\section{Implications and Requirements for Future Work}\label{sec:implications}

\subsection{Status of the phenomenological framework}

The phenomenological framework of Ref.~\cite{Golden2026}---the scaling ansatz
$\rho_\Lambda = \Xi\,\MPl^4$ with $w=-1$ as an assumption---is unaffected by
these closures. Its cosmological consistency checks (MCMC fits to
Planck/BAO/SN data), fine-tuning reduction ($10^{120} \to {\sim}10^5$), and
qualitative parity/birefringence consistency survive as phenomenological
observations. The framework was designed to be robust under exactly this
outcome.

However, the closures do establish the framework's ceiling: it is a
\emph{motivated phenomenological model}, not a first-principles derivation
of dark energy, and its minimal model class lacks both a derivation of
$w=-1$ and a distinctive observational signature separating it from
$\Lambda$CDM$+$ALP.

\subsection{Requirements for next-generation models}

The structural lessons from Sec.~\ref{sec:combined} translate into four
concrete requirements for any candidate next-generation geometric
dark-energy model:

\begin{description}[style=nextline]
  \item[R1: Must survive reduction without washing out.]
    The mechanism must not depend on degrees of freedom that are eliminated
    by algebraic constraint equations. If torsion is algebraic, nothing that
    relies on it persists in the reduced action.

  \item[R2: Must address tiny-scale naturalness.]
    The model must explain why the relevant scale is $\sim 10^{-122}\,\MPl^4$
    through a structural mechanism (symmetry, sequestering, relaxation), not
    merely by setting a free parameter to the observed value.

  \item[R3: Must offer a distinctive prediction.]
    The model must produce at least one observable or mechanism that is not
    reproducible by $\Lambda$CDM plus a generic axion-like particle.

  \item[R4: Must fail cleanly if wrong.]
    The model must be testable with predefined gates and kill criteria,
    producing documentable negative results if it fails.
\end{description}

\subsection{Candidate directions beyond the minimal model}

Three directions satisfy at least one of these requirements while addressing
at least one structural lesson:
\begin{itemize}
  \item \emph{Propagating torsion} (Poincar\'e gauge theory): torsion becomes
    dynamical, avoiding the algebraic wash-out (R1). Risk: ghost
    instabilities in the massive torsion sector.
  \item \emph{Symmetry-protected geometric pseudoscalars:} a pseudoscalar
    with a protected mass and non-topological geometric coupling, avoiding
    ALP collapse (R3) and addressing naturalness (R2). Risk: may reduce
    to a standard ALP in disguise.
  \item \emph{Vacuum sequestering:} global constraints enforce radiative
    stability of $\Lambda$, addressing naturalness directly (R2). Risk:
    Weinberg no-go theorem; potentially no distinctive observable (R3).
\end{itemize}
Each requires its own canonical problem statement and gate structure before
computation.


% ------------------------------------------------------------------
\section{Reproducibility and Repository Layout}\label{sec:reproducibility}

Track~B and Branch~G v1 computations are symbolic or semi-analytic,
implemented in Python~3.14 with SymPy~1.14, NumPy, SciPy, and Matplotlib.
No GPU or specialized hardware is required. Routes~T1 and S1 are
literature-review and structural-assessment closures with no original
computation; their documentation is archived in the repository. The full
repository is organized as follows.

\subsection{Frozen documents (not to be edited)}

\begin{longtable}{@{}p{0.42\textwidth}p{0.52\textwidth}@{}}
\toprule
File & Content \\
\midrule
\texttt{06\_canonical\_problem\_statement.md} &
  Track~B canonical problem statement (v3, frozen).
  Defines 7 questions, 3 gates, order parameter, symmetry logic,
  regulator strategy. \\
\texttt{06a\_frozen\_assumptions\_and\_change\_log.md} &
  Track~B freeze log with reopening criteria and v1$\to$v2$\to$v3
  change history. \\
\texttt{07\_track\_b\_closure.md} &
  Track~B closure document. Four load-bearing negative results,
  artifact inventory, impact assessment. \\
\texttt{08\_branchG\_canonical\_problem\_statement.md} &
  Branch~G v1 canonical problem statement (v2, frozen).
  6 questions, 3 gates, FM-G6 early kill, scope control. \\
\texttt{08a\_branchG\_freeze\_log.md} &
  Branch~G v1 freeze log with Approach~A/B label clarification
  and Track~B comparison table. \\
\bottomrule
\end{longtable}

\subsection{Computation scripts}

\begin{longtable}{@{}p{0.40\textwidth}p{0.08\textwidth}p{0.42\textwidth}@{}}
\toprule
Script & Type & Content \\
\midrule
\texttt{derivations/comp1\_torsion\_elimination.py} & Symbolic &
  Torsion elimination in ECH gravity. Derives perfect-square
  structure~\eqref{eq:Sfour}, coupling constants, $\BI\to\infty$
  limit check. \\
\texttt{derivations/comp2\_fierz\_rearrangement.py} & Symbolic &
  Full $16\times 16$ Fierz rearrangement using explicit gamma
  matrices in the Dirac representation. Channel decomposition. \\
\texttt{derivations/comp2b\_fierz\_verification.py} & Symbolic &
  Independent verification of Fierz results. VA vanishing proof,
  sign analysis, Fierz completeness check. \\
\texttt{derivations/comp3\_one\_loop\_veff.py} & Semi-analytic &
  One-loop $\Veff$ with curvature corrections. $\BI$-scan for
  nontrivial minima (none found). Multi-flavor analysis.
  Generates Figs.~\ref{fig:GSP}--\ref{fig:veff}. \\
\texttt{derivations/comp3b\_gap\_equation\_check.py} & Numerical &
  NJL gap equation in proper-time regularization. Critical
  coupling comparison ($\Gsp/\Gcrit\approx 0.006$). Curvature
  catalysis scan. \\
\texttt{derivations/branchG\_phase1\_operator\_analysis.py} & Analytic &
  Branch~G Phase~1: three-channel analysis of whether $\BI$
  enters the strict one-loop Dirac operator. \\
\bottomrule
\end{longtable}

\subsection{Analysis notes and status reports}

\begin{longtable}{@{}p{0.42\textwidth}p{0.52\textwidth}@{}}
\toprule
File & Content \\
\midrule
\texttt{notes/gate1\_status\_20260313.md} &
  Gate~1 failure report: three results that kill the S/P condensate
  route. Remaining active branches. \\
\texttt{notes/comp3\_results\_analysis.md} &
  Fierz result analysis: VA cross-term structure, multi-flavor
  implications, Possibility A vs B resolution. \\
\texttt{notes/branchG\_phase1\_closure.md} &
  Branch~G Phase~1 closure: FM-G6 declaration, three-channel
  verdict, literature comparison. \\
\bottomrule
\end{longtable}

\subsection{Route~T1 and S1 documentation}

Routes~T1 and S1 were closed through literature review and structural
assessment. No original symbolic computation was required. The following
documents record the assessments:

\begin{longtable}{@{}p{0.42\textwidth}p{0.52\textwidth}@{}}
\toprule
File & Content \\
\midrule
\texttt{next\_branches/T1\_dynamical\_immirzi/\newline
01\_canonical\_problem\_statement.md} &
  T1 canonical problem statement. 6 questions, 3 gates, 6 failure modes,
  4-phase computation sequence. \\
\texttt{next\_branches/T1\_dynamical\_immirzi/\newline
06\_phase1\_literature\_review\_questions.md} &
  T1 Phase~1 checklist: 6 sharp questions with kill criteria, 7-day reading
  plan, assessment memo template. \\
\texttt{next\_branches/S1\_parity\_cmb/\newline
01\_canonical\_problem\_statement.md} &
  S1 canonical problem statement. 4 target observables, 3 gates, 5 failure
  modes. \\
\texttt{next\_branches/S1\_parity\_cmb/\newline
05\_phase1\_assessment\_checklist.md} &
  S1 Phase~1 checklist: 4 questions with kill criteria, 5-day timeline,
  mapping memo template. \\
\texttt{next\_branches/\newline
01\_sequencing\_and\_decision\_gates.md} &
  Master sequencing table, gate dependencies, critical path analysis for
  all four branches (T1, T2, S1, S2). \\
\bottomrule
\end{longtable}

\subsection{Figures}

\begin{longtable}{@{}p{0.42\textwidth}p{0.52\textwidth}@{}}
\toprule
File & Content \\
\midrule
\texttt{figures/G\_SP\_vs\_gamma.png} &
  $\Gsp$ vs.\ $\BI$ plot (Fig.~\ref{fig:GSP}).
  Generated by \texttt{comp3\_one\_loop\_veff.py}. \\
\texttt{figures/veff\_repulsive\_vs\_attractive.png} &
  $\Veff$ comparison plot (Fig.~\ref{fig:veff}).
  Generated by \texttt{comp3\_one\_loop\_veff.py}. \\
\bottomrule
\end{longtable}

\subsection{Reproduction instructions}

To reproduce all results:
\begin{enumerate}
  \item Install Python~3.14 with \texttt{pip install sympy numpy scipy
    matplotlib}.
  \item Navigate to \texttt{ir\_vacuum\_program/derivations/}.
  \item Run scripts in order:
    \texttt{comp1\_torsion\_elimination.py} $\to$
    \texttt{comp2\_fierz\_rearrangement.py} $\to$
    \texttt{comp2b\_fierz\_verification.py} $\to$
    \texttt{comp3\_one\_loop\_veff.py} $\to$
    \texttt{comp3b\_gap\_equation\_check.py} $\to$
    \texttt{branchG\_phase1\_operator\_analysis.py}.
  \item All output is printed to stdout. Figures are written to
    \texttt{../figures/}.
  \item No script depends on external data, GPU, or network access.
\end{enumerate}

\subsection{What is symbolic vs.\ numerical}

\begin{longtable}{@{}p{0.30\textwidth}p{0.65\textwidth}@{}}
\toprule
Category & Scripts \\
\midrule
Purely symbolic (SymPy) &
  \texttt{comp1}, \texttt{comp2}, \texttt{comp2b}: gamma matrix algebra,
  Fierz traces, coupling constants. Exact algebraic results. \\
Semi-analytic &
  \texttt{comp3}: symbolic setup + numerical $\BI$-scan of $\Veff$.
  \texttt{comp3b}: analytic gap equation + numerical root-finding
  and catalysis estimates. \\
Structural/analytic &
  \texttt{branchG\_phase1}: logical argument verified against literature;
  no numerics beyond illustrative coefficient values. \\
\bottomrule
\end{longtable}


% ------------------------------------------------------------------
\section{Conclusion}\label{sec:conclusion}

Four routes to deriving the late-time vacuum term or identifying distinctive
observables in the spin-torsion dark-energy framework were tested and found
to fail:
\begin{enumerate}
  \item The NJL condensate route (Track~B) fails because the
    scalar/pseudoscalar channel is repulsive at the LQG Barbero-Immirzi value
    $\BI=0.274$, the gravitational four-fermion coupling is ${\sim}175\times$
    subcritical, and gravitational catalysis is exponentially suppressed.
  \item The strict one-loop effective-action route (Branch~G v1) fails because
    the fermion determinant after torsion elimination is the standard
    $\BI$-independent curved-space Dirac determinant, with no novel Holst
    content at this order.
  \item The dynamical Immirzi field extension (Route~T1) fails because torsion
    elimination reduces the dynamical $\theta(x)$ field to a standard
    axion-like particle with no novel gravitational content, no mechanism for
    generating $V(\theta)\sim\Lambda$, and cosmological dynamics yielding
    $w=+1$ (stiff matter).
  \item The parity-sensitive CMB phenomenology assessment (Route~S1) fails
    because the minimal framework has no photon coupling, the assumed coupling
    produces constant-$\beta$ birefringence identical to generic ALP
    phenomenology, and the parameterization is unfalsifiable.
\end{enumerate}

These closures share a common structural origin: in the minimal ECH$+$Dirac
model, torsion elimination removes all Holst-sector structure, and the
Nieh-Yan topological density carries no local dynamical content that survives
this elimination. The resulting physics is standard GR~$+$~fermions~$+$~a
constrained four-fermion interaction, with no mechanism for late-time vacuum
energy and no distinctive observational signature.

These are clean negative results with a unifying structural explanation:
algebraic torsion in the minimal ECH model washes out completely after
elimination, and the Nieh-Yan density's topological nature in four dimensions
ensures that even the simplest structural extension (dynamical Immirzi field)
collapses to standard physics.

The value of these results is threefold. First, they map the landscape of
what does not work within the most natural minimal model, preventing
duplication of effort. Second, they identify the precise structural
bottleneck---algebraic torsion and topological Nieh-Yan---that any successor
model must overcome. Third, they establish a reproducible methodology
(canonical problem statements, predeclared gates, frozen scope) for
disciplined model testing that can be applied to future candidates.

The phenomenological framework of Ref.~\cite{Golden2026} survives as a
motivated ansatz; the first-principles derivation of $w=-1$ and the
identification of a distinctive observable remain open problems requiring
physics beyond the minimal model class.


% ------------------------------------------------------------------
% REFERENCES
% ------------------------------------------------------------------
\begin{thebibliography}{99}

\bibitem{Golden2026}
H.~Golden,
``Geometric Dark Energy from Spin-Torsion Cosmology: Phenomenological
Constraints and Correlated Signatures,''
Preprint HUBIFY-2026-001 (2026).

\bibitem{Freidel2005}
L.~Freidel, D.~Minic, and T.~Takeuchi,
``Quantum gravity, torsion, parity violation and all that,''
Phys.\ Rev.\ D~\textbf{72}, 104002 (2005)
[arXiv:hep-th/0507253].

\bibitem{ShapiroTeixeira2014}
I.~L.~Shapiro and P.~M.~Teixeira,
``Quantum Einstein-Cartan theory with the Holst term,''
Class.\ Quant.\ Grav.\ \textbf{31}, 185002 (2014)
[arXiv:1402.4854].

\bibitem{Chattopadhyay2023}
P.~Chattopadhyay,
``One-loop effective action in chiral Einstein-Cartan gravity,''
Int.\ J.\ Mod.\ Phys.\ A~\textbf{38}, 2350186 (2023)
[arXiv:2310.10405].

\bibitem{GorbarGusynin2008}
E.~V.~Gorbar and V.~P.~Gusynin,
``Dynamical chiral symmetry breaking on a brane in reduced QED,''
Phys.\ Rev.\ D~\textbf{73}, 105028 (2006).
% Note: representative reference for gravitational catalysis framework.

\bibitem{Vassilevich2003}
D.~V.~Vassilevich,
``Heat kernel expansion: User's manual,''
Phys.\ Rept.\ \textbf{388}, 279 (2003)
[arXiv:hep-th/0306138].

\bibitem{Shapiro2002}
I.~L.~Shapiro,
``Physical aspects of the space-time torsion,''
Phys.\ Rept.\ \textbf{357}, 113 (2002)
[arXiv:hep-th/0103093].

\bibitem{Mercuri2009}
R.~Mercuri,
``Peccei-Quinn mechanism and the Barbero-Immirzi parameter,''
Phys.\ Rev.\ Lett.\ \textbf{103}, 081302 (2009)
[arXiv:0902.2764].

\bibitem{CalcagniMercuri2009}
G.~Calcagni and R.~Mercuri,
``The Barbero-Immirzi field and its cosmological implications,''
Phys.\ Rev.\ D~\textbf{79}, 104003 (2009)
[arXiv:0902.0957].

\bibitem{TaverasYunes2008}
V.~Taveras and N.~Yunes,
``The Barbero-Immirzi parameter as a scalar field:
$K$-inflation from loop quantum gravity?,''
Phys.\ Rev.\ D~\textbf{78}, 064070 (2008)
[arXiv:0812.3572].

\bibitem{TorresGomez2009}
A.~Torres-Gomez, K.~Krasnov, and C.~Scarinci,
``Barbero-Immirzi parameter as a scalar field: effects of fermions,''
Phys.\ Rev.\ D~\textbf{80}, 104029 (2009)
[arXiv:0904.4877].

\bibitem{LattanziMercuri2010}
M.~Lattanzi and R.~Mercuri,
``A solution of the strong CP problem via the Peccei-Quinn mechanism
through the Nieh-Yan modified gravity and cosmological implications,''
Phys.\ Rev.\ D~\textbf{81}, 125015 (2010)
[arXiv:0911.2698].

\bibitem{Magueijo2013}
J.~Magueijo, T.~G.~Zlosnik, and T.~W.~B.~Kibble,
``Cosmology with a spin,''
Phys.\ Rev.\ D~\textbf{87}, 063504 (2013)
[arXiv:1212.0585].

\bibitem{EskiltKomatsu2022}
J.~R.~Eskilt and E.~Komatsu,
``Improved constraints on cosmic birefringence from the WMAP and
Planck cosmic microwave background polarization data,''
Phys.\ Rev.\ D~\textbf{106}, 063503 (2022)
[arXiv:2205.13962].

\bibitem{DiegoPalazuelos2025}
P.~Diego-Palazuelos and E.~Komatsu,
``Cosmic birefringence from the Atacama Cosmology Telescope DR6,''
arXiv:2501.10850 (2025).

\bibitem{canonical_v3}
Canonical problem statement for Track~B, version~3 (frozen 2026-03-13).
Repository file: \texttt{06\_canonical\_problem\_statement.md}.

\bibitem{branchG_canonical}
Canonical problem statement for Branch~G v1, version~2 (frozen 2026-03-13).
Repository file: \texttt{08\_branchG\_canonical\_problem\_statement.md}.

\bibitem{comp1}
Computation~1: Torsion elimination.
Repository file: \texttt{derivations/comp1\_torsion\_elimination.py}.

\bibitem{comp2}
Computation~2: Fierz rearrangement.
Repository file: \texttt{derivations/comp2\_fierz\_rearrangement.py}.

\bibitem{comp2b}
Computation~2b: Fierz verification.
Repository file: \texttt{derivations/comp2b\_fierz\_verification.py}.

\bibitem{comp3}
Computation~3: One-loop effective potential.
Repository file: \texttt{derivations/comp3\_one\_loop\_veff.py}.

\bibitem{comp3b}
Computation~3b: Gap equation check.
Repository file: \texttt{derivations/comp3b\_gap\_equation\_check.py}.

\bibitem{branchG_analysis}
Branch~G Phase~1 analysis.
Repository file: \texttt{derivations/branchG\_phase1\_operator\_analysis.py}.

\end{thebibliography}

\end{document}
